% ============================================================
% rq2_cognition.tex — Fluid Intelligence × Cortical Surface Area (HCP)
% ============================================================

\section{Fluid Intelligence and Cortical Surface Area in Young Adults}
\label{sec:rq2}

\subsection{Research Question and Hypothesis}

The neural basis of human intelligence has been studied extensively, with the Parieto-Frontal Integration Theory (P-FIT) providing the most comprehensive neuroanatomical framework to date~\cite{jung2007pfit}.
P-FIT posits that intelligent behavior depends on the efficient integration of information across a distributed network anchored in the dorsolateral prefrontal cortex and the superior and inferior parietal lobules, with additional contributions from temporal and occipital association areas~\cite{jung2007pfit, colom2009human}.
While many studies have examined the relationship between intelligence and gray matter \textit{volume}~\cite{pietschnig2015meta}, volume confounds two geometrically distinct cortical properties: thickness and surface area.
Surface area, which reflects the degree of cortical folding and the tangential expansion of cortical columns, has a substantially different genetic architecture from thickness~\cite{grasby2020genetic} and may therefore relate to cognitive ability through distinct biological mechanisms.
Few studies have isolated the surface-area component in a sample large enough to test region-wise associations with adequate statistical power.

We therefore asked: \textit{Is fluid intelligence associated with regional cortical surface area in healthy young adults?}

\textbf{Hypothesis.}
Drawing on the P-FIT model~\cite{jung2007pfit} and the observation that cortical expansion is greatest in association cortices~\cite{grasby2020genetic}, we hypothesized that higher fluid intelligence scores would be associated with greater cortical surface area, with the strongest associations in parietal, frontal, and temporal regions.

\subsection{Results}

Pearson correlations were computed between PMAT24 scores and cortical surface area for each of the 68 DK ROIs.
After Benjamini-Hochberg FDR correction at \textit{q}\,<\,.05, all 68 regions showed a statistically significant positive correlation between fluid intelligence and surface area.
The universality of this effect was striking: even the weakest regional correlation reached statistical significance after correction.

The five strongest associations were concentrated in the temporal lobe:
\begin{itemize}[nosep]
  \item Right inferior temporal: \textit{r}(1{,}104)\,=\,.290, \textit{p}\,<\,.001, 95\% CI [.234, .344]
  \item Left middle temporal: \textit{r}(1{,}104)\,=\,.278, \textit{p}\,<\,.001, 95\% CI [.221, .333]
  \item Right middle temporal: \textit{r}(1{,}104)\,=\,.273, \textit{p}\,<\,.001, 95\% CI [.216, .328]
  \item Left inferior temporal: \textit{r}(1{,}104)\,=\,.265, \textit{p}\,<\,.001, 95\% CI [.208, .320]
  \item Right superior temporal: \textit{r}(1{,}104)\,=\,.258, \textit{p}\,<\,.001, 95\% CI [.201, .314]
\end{itemize}
Frontal and parietal regions also showed robust positive correlations, with \textit{r}-values typically between .15 and .25.
The overall pattern is visualized in Figure~\ref{fig:rq2-scatter}.

\begin{figure}[htbp]
  \centering
  \includegraphics[width=0.95\textwidth]{hcp/cognition_scatter_top5.pdf}
  \caption{Scatter plots of fluid intelligence (PMAT24 score) versus cortical surface area for the five regions with the strongest Pearson correlations.
           Each point represents one participant (\textit{N}\,=\,1{,}104).
           All 68 regions reached FDR-corrected significance.}
  \label{fig:rq2-scatter}
\end{figure}

\subsection{Discussion}

The universal positive association between fluid intelligence and cortical surface area across all 68 DK regions is consistent with the view that intelligence is not localized to a handful of ``smart'' regions but instead reflects the global efficiency of a broadly distributed cortical network~\cite{deary2010neuroscience}.
This finding extends prior work on the P-FIT model~\cite{jung2007pfit}, which emphasized parietal and frontal nodes, by demonstrating that temporal association areas show the strongest surface-area associations with fluid reasoning.
The inferior and middle temporal gyri are involved in visual object recognition, semantic processing, and the integration of bottom-up perceptual evidence with top-down expectations, processes that are central to the matrix-reasoning task used to assess fluid intelligence in the HCP~\cite{colom2009human}.

Several features of our results merit emphasis.
First, the largest correlations (\textit{r}\,$\approx$\,.29) are modest by conventional standards but are among the largest region-level associations reported in single-cohort studies of comparable size~\cite{pietschnig2015meta}.
The magnitude is consistent with the meta-analytic estimate of \textit{r}\,$\approx$\,.25 for the overall association between total brain volume and intelligence~\cite{pietschnig2015meta}, suggesting that surface area captures much of the structural variance that is relevant to cognitive ability.
Second, the specificity to surface area, rather than thickness, is theoretically informative.
Cortical surface area is largely determined by the number of cortical columns, which is established prenatally and shows high heritability~\cite{grasby2020genetic}.
Cortical thickness, by contrast, reflects the laminar organization of individual columns and is more sensitive to environmental and experiential factors.
The stronger association with surface area therefore suggests that the structural substrate of fluid intelligence in young adults is primarily a reflection of early cortical development rather than ongoing experience-dependent plasticity.

A caveat is that the PMAT24 is a relatively brief 24-item measure that assesses only nonverbal matrix reasoning, which captures one facet of fluid intelligence.
Studies using more comprehensive cognitive batteries may reveal a more differentiated regional pattern.
